\documentclass[11pt]{article}
\usepackage[margin=1in]{geometry}
\usepackage{amsmath, amssymb, amsthm}
\usepackage{booktabs}
\usepackage{hyperref}
\usepackage{xcolor}
\usepackage{listings}
\usepackage{xurl}

\hypersetup{colorlinks=true, urlcolor=blue, linkcolor=black, citecolor=black}

\lstset{
  basicstyle=\ttfamily\small,
  numbers=left,
  numberstyle=\tiny,
  frame=single,
  breaklines=true,
  showstringspaces=false,
  language=Python,
  keywordstyle=\color{blue},
  commentstyle=\color{gray}\itshape,
  stringstyle=\color{teal},
}

\title{Reconstruction Is Not Replication: \\
       Cross-Validating Sparse Autoencoder Features on Pythia-160M}

\author{Aung Maw\\
        Skyline College, San Francisco\\
        \texttt{irving46764@gmail.com} \\
        \href{https://github.com/OE-GOD}{\texttt{github.com/OE-GOD}}}

\date{May 2026}

\begin{document}
\maketitle

\begin{abstract}
We train a sparse autoencoder (SAE) on the layer-6 residual stream of
Pythia-160M and evaluate it along six independent quality axes, including
three cross-validation tests that the field rarely reports together:
cross-width consistency (training the same SAE at $N \in \{4096, 16384,
65536\}$ features), cross-seed replication (training at the same
hyperparameters with a different random seed), and cross-architecture
agreement (training a Matryoshka SAE \cite{bussmann2025matryoshka} with
nested-$k$ bottlenecks).
The reference SAE achieves $98.4\%$ variance explained, $95.4\%$ loss
recovered under cross-entropy substitution (a competitive number for the
small-model band), and $88.5\%$ auto-interp monosemanticity rate over the
$26$ most-active features.
The same SAE has only $19.9\%$ of features at cosine similarity $> 0.9$
across a different random seed, and only $10.6\%$ at the same threshold
across the Matryoshka architecture. Combining all four cross-validation
conditions, only $2.14\%$ of features ($351$ of $16{,}384$) survive at
cos $> 0.9$. Of the $26$ features that an automatic interpreter labeled
as monosemantic, only $3$ are fully stable: features encoding BibTeX/LaTeX
citation references, file paths, and logical operators / negation.
We argue that the standard suite of SAE quality metrics
(reconstruction, sparsity, monosemanticity rate) systematically
\emph{overstates} feature reliability, by approximately $8\times$ on this
setup. We release a multi-metric scorecard tool and a per-feature
stability navigator that surface this gap on any SAE.\footnote{Code:
\url{https://github.com/OE-GOD/sae-pythia-160m}.}
\end{abstract}

\section{Introduction}

Sparse autoencoders (SAEs) have become the dominant tool for converting a
transformer's internal activations into a sparse, interpretable basis
\cite{bricken2023monosemanticity, templeton2024scaling}.
The standard quality story has three parts: reconstruction (variance
explained or MSE), sparsity (an $L_0$ target enforced by ReLU+$L_1$
\cite{bricken2023monosemanticity}, top-$k$ \cite{gao2024topk}, or related
mechanisms \cite{rajamanoharan2024gated, rajamanoharan2024jumprelu}), and
interpretability (the fraction of recovered features that a human or LLM
inspector labels as monosemantic). Most published SAE work reports these
three.

This paper asks a question the standard suite cannot answer: if you train
the same SAE again --- with a different random seed, a different width,
or a slightly different architectural objective --- how many of the
features it recovered are the same?

Our finding: \emph{very few}. On Pythia-160M layer 6 with $N=16{,}384$
features and $k=64$:
\begin{itemize}
\item Variance explained is $98.4\%$; CE-delta loss recovered is $95.4\%$.
\item Auto-interp on the top $26$ features classifies $88.5\%$ as
monosemantic.
\item But cross-seed cosine similarity is $> 0.9$ for only $19.9\%$ of
features; cross-Matryoshka is $> 0.9$ for only $10.6\%$.
\item Of the auto-labeled features, only $3$ of $26$ are fully stable
across width, seed, and architecture. The rest, including all
$11$ apparent variants of ``newline tokens'' that the SAE found, collapse
under cross-validation.
\end{itemize}

These observations have two implications. First, when the field reports
SAE quality as a single reconstruction or interpretability number, that
number can overstate the reliability of recovered features by nearly an
order of magnitude. Second, the small subset of features that does
survive cross-validation is, plausibly, the right basis for downstream
work --- circuit discovery, causal interventions, agent behavior auditing
--- because those features are not artifacts of the particular SAE
training procedure.

\paragraph{Contributions.}
\begin{enumerate}
\item Empirical: a six-axis evaluation of a TopK SAE on Pythia-160M layer
6, including cross-width (\S\ref{sec:width}), cross-seed
(\S\ref{sec:seed}), and cross-architecture (\S\ref{sec:matryoshka})
replication results. The cross-validation numbers are roughly
$\sim 5\times$--$10\times$ worse than the reconstruction and
interpretability numbers on the same SAE.
\item A multi-metric scorecard tool that reports all six axes side-by-side
(\S\ref{sec:scorecard}). On our reference SAE the strongest axis is
$0.984$ and the weakest is $0.106$, a $\sim 88$ percentage-point gap.
\item A per-feature stability navigator that reports a feature's
best-match identity in every comparison SAE
(\S\ref{sec:navigator}). The output is a shortlist of $351$ fully stable
features ($2.14\%$ of the total) that we treat as the strongest ``real
feature'' candidates in this setup.
\end{enumerate}

\paragraph{Scope and honesty.}
This is a single-model, single-layer study on a small open transformer
($160$M params) with $1\text{M}$ training tokens and a $\$50$ compute
budget. We make no claims about frontier-scale models. We do claim that
on this benchmark the gap between reconstruction-style and
replication-style metrics is large enough to deserve attention in any
SAE quality report.

\section{Background}

The SAE setup is standard. For an activation $h \in \mathbb{R}^d$ at a
target layer,
\begin{align}
f &= \sigma(W_{\text{enc}} (h - b_{\text{dec}}) + b_{\text{enc}}), \quad
f \in \mathbb{R}^N, \\
\hat{h} &= W_{\text{dec}} f + b_{\text{dec}},
\end{align}
with $\sigma$ enforcing sparsity. We use \textsc{TopK}
\cite{gao2024topk}, which keeps the $k$ largest pre-activations and zeros
the rest. This gives $L_0 = k$ by construction and removes the
$L_1$-shrinkage problem \cite{rajamanoharan2024gated} we observed in an
initial vanilla-$L_1$ run (Appendix~\ref{app:l1-shrinkage}).

The Matryoshka SAE variant \cite{bussmann2025matryoshka} introduces
nested-$k$ reconstruction bottlenecks: for $k_1 < k_2 < \dots < k_M$,
features active at smaller $k_i$ are a strict subset of features at larger
$k_j$, and the loss sums reconstruction over all levels:
\begin{equation}
\mathcal{L}_{\text{matryoshka}} = \sum_{i=1}^{M} \| h - W_{\text{dec}} f^{(k_i)}_{\text{topk}} - b_{\text{dec}} \|_2^2.
\end{equation}
We use this variant for the cross-architecture comparison in
\S\ref{sec:matryoshka}.

\section{Setup}

\paragraph{Model and data.}
Pythia-160M \cite{biderman2023pythia} ($d_{\text{model}}=768$, $12$
layers), accessed via TransformerLens. Target: residual stream at the
output of layer $6$, hook \texttt{blocks.6.hook\_resid\_post}.
Training corpus: \texttt{NeelNanda/pile-10k}, a $10\text{k}$-document
subset of the Pile. We tokenise the first $\sim 8{,}000$ documents into a
flat stream and reshape into $7{,}812$ sequences of length $128$, giving
$999{,}936$ training tokens. Held-out evaluation tokens are drawn from
documents indexed $\geq 8{,}001$.

\paragraph{SAE training.}
The reference SAE (referred to as the ``$16$k SAE'') has $N = 16{,}384$
features ($\sim 21 \times$ expansion over $d_{\text{model}}$), $k = 64$
(TopK), $b_{\text{dec}}$ initialised to the empirical activation mean,
tied initialisation ($W_{\text{enc}} = W_{\text{dec}}^\top$), decoder
column unit-norm constraint enforced after each Adam step, $\text{lr} =
10^{-3}$, batch size $4{,}096$, $20{,}000$ training steps.
Three additional SAEs are trained for cross-validation:
\begin{itemize}
\item $N = 4{,}096$ (cross-width, smaller, \S\ref{sec:width})
\item $N = 65{,}536$ (cross-width, larger, \S\ref{sec:width})
\item Second random seed (\texttt{torch.manual\_seed(42)}, all other
hyperparameters identical, \S\ref{sec:seed})
\item Matryoshka SAE at $k_{\text{levels}} = \{16, 64, 256\}$
(\S\ref{sec:matryoshka})
\end{itemize}

\paragraph{Hardware and compute.}
Apple Silicon (M-series) via PyTorch MPS. All training and evaluation runs
fit within a $\$50$ compute budget; in practice the project consumed zero
paid compute (all on local hardware over $\sim 24$ hours of
background time).

\section{Single-SAE Results (Baseline)}

\subsection{Reconstruction and sparsity}

On $50{,}000$ held-out activations from the training-document set:
variance explained $98.43\%$, mean $L_0 = 63.9 \pm 1.1$ (target $k=64$,
exact), reconstruction MSE summed over $d_{\text{model}}$ is $25.74$.
Dead features (features that never activate): $855 / 16{,}384 = 5.22\%$.

\subsection{CE-delta evaluation}
\label{sec:ce-delta}

Following \cite{lieberum2024gemma}, we evaluate the SAE by substituting
its reconstruction into the language model's forward pass at the trained
layer, and measuring the increase in downstream cross-entropy on
held-out tokens. For $100{,}000$ held-out tokens:
\begin{align}
\text{CE}_{\text{clean}} &= 3.107 \text{ nats/token} \\
\text{CE}_{\text{SAE-patched}} &= 3.383 \text{ nats/token} \\
\text{CE}_{\text{zero-ablate (mean)}} &= 9.122 \text{ nats/token}
\end{align}
The SAE-patched run substitutes $\hat{h} = \text{SAE}(h)$ for $h$ at layer
$6$; the zero-ablation run substitutes the data mean
($b_{\text{dec}}$). The loss-recovered metric is
$1 - \Delta\text{CE} / (\text{CE}_{\text{ablate}} - \text{CE}_{\text{clean}})
= 95.4\%$, in the strong end of the published small-model band
\cite{bricken2023monosemanticity, lieberum2024gemma}.

\subsection{Auto-interpretation catalog}
\label{sec:autointerp}

We extract top-$10$ max-activating snippets per feature on the $1\text{M}$-
token training stream, then submit the top $26$ features (ranked by a
mixture of peak and mean-of-top-$K$ activation) to Kimi-K2
(\texttt{moonshot-v1-32k}) for auto-interpretation. Each receives a
one-sentence concept label and a classification of MONOSEMANTIC,
POLYSEMANTIC, or UNCLEAR. Total API cost: $\$0.02$.

Of the $26$:
\begin{itemize}
\item MONOSEMANTIC: $23$ ($88.5\%$).
\item POLYSEMANTIC: $3$ ($11.5\%$).
\item UNCLEAR: $0$.
\end{itemize}

Representative monosemantic features (full table in
Appendix~\ref{app:features}): \emph{Decimal points in numerical contexts}
(f1989), \emph{Exponentiation notation in mathematical contexts} (f5196),
\emph{Tokens related to file paths and directory structures} (f5747),
\emph{Newline tokens inside code blocks} (f7448),
\emph{BibTeX/LaTeX citation references} (f12117),
\emph{Logical operators and negation symbols} (f12697).

\textbf{Apparent feature splitting.} Approximately $11$ of the $23$
monosemantic features received label variants of ``newline tokens'':
\emph{newline tokens in various contexts}, \emph{newline tokens inside
XML/HTML/CSS blocks}, \emph{newline tokens inside code blocks}, and so on.
The SAE has clearly split a single semantic category into many features.
The labeling LLM identifies each as monosemantic but cannot give distinct
labels because the splits are subtle. This is direct evidence for
feature splitting at $N = 16{,}384$ on this model.

\section{Cross-Validation Results (the New Material)}

\subsection{Cross-width sweep}
\label{sec:width}

Three SAEs trained at $N \in \{4{,}096, 16{,}384, 65{,}536\}$ with all
other hyperparameters held fixed. Results in
Table~\ref{tab:pareto}.

\begin{table}[h]
\centering
\caption{Cross-width Pareto. Wider SAEs do not buy meaningful
reconstruction beyond $N \approx 4\text{k}$ on this data scale, and
dead-feature counts explode at $N=65{,}536$.}
\label{tab:pareto}
\begin{tabular}{lcccc}
\toprule
Width $N$ & MSE & Var.\ explained & Dead features & Alive features \\
\midrule
$4{,}096$  & $33.11$ & $0.9782$ & $1 \,(0.02\%)$  & $4{,}095$ \\
$16{,}384$ & $25.74$ & $0.9843$ & $855 \,(5.22\%)$ & $15{,}529$ \\
$65{,}536$ & $23.68$ & $0.9853$ & $41{,}698 \,(63.63\%)$ & $23{,}838$ \\
\bottomrule
\end{tabular}
\end{table}

Two observations. \textbf{The Pareto curve in width is nearly flat above}
$N = 4{,}096$: going $4\text{k} \to 64\text{k}$ ($16\times$ more
capacity) buys $0.71$ percentage points of variance explained.
\textbf{Dead features explode at $N = 65{,}536$}: $63.6\%$ of features
never fire on any input.

\textbf{Feature merging at small width.} We find one $4\text{k}$ feature
(f2175) is the best cosine-similarity match for five different $16\text{k}$
newline-variant features (cos $0.34$ to $0.39$). Smaller widths collapse
multiple features that nominally differ at larger widths.

\textbf{Feature splitting at large width does not survive seed change.}
Anticipating \S\ref{sec:seed}: most of the apparent $16\text{k}$ features
that look like ``finer-grained splits'' of $4\text{k}$ features do not
have stable counterparts in a second random-seed SAE.

\subsection{Cross-seed replication}
\label{sec:seed}

We train a second SAE at identical hyperparameters with
\texttt{torch.manual\_seed(42)} (the reference SAE used the default
implicit seed). The two SAEs have statistically indistinguishable summary
metrics: variance explained $98.43\%$ vs.\ $98.44\%$, dead features
$5.22\%$ vs.\ $5.47\%$, $L_0 = 63.9$ for both.

For each feature in the reference, we find its best-match decoder column
in the second SAE under cosine similarity (with decoder columns unit-norm
across both). Distribution of best-match cosine similarities:

\begin{center}
\begin{tabular}{lcc}
\toprule
Threshold & Count above & Fraction \\
\midrule
cos $> 0.50$ & $12{,}892$ & $0.787$ \\
cos $> 0.70$ & $9{,}783$ & $0.597$ \\
cos $> 0.90$ & $3{,}254$ & $\mathbf{0.199}$ \\
cos $> 0.95$ & $952$ & $0.058$ \\
cos $> 0.99$ & $14$ & $0.001$ \\
\bottomrule
\end{tabular}
\end{center}

\textbf{Only $19.9\%$ of features have a cross-seed match above
$\text{cos} = 0.9$.} Same model, same data, same hyperparameters, same
architecture: only $1$ in $5$ features replicates.

\textbf{All $11$ ``newline-variant'' features in the reference SAE collapse
to a single feature in the seed-$42$ SAE.} Specifically, the
best-match feature in the seed-$42$ SAE for each of the labeled newline
features is the same feature (f12024) with cosine similarities of
$0.34$--$0.46$. This directly shows that the apparent ``feature splitting''
observed at $N = 16{,}384$ is substantially a training-noise artifact
rather than the SAE discovering finer-grained concepts.

\subsection{Cross-architecture: Matryoshka vs.\ flat TopK}
\label{sec:matryoshka}

We train a Matryoshka SAE \cite{bussmann2025matryoshka} on the same
$1\text{M}$ activations with $k_{\text{levels}} = \{16, 64, 256\}$ and
the same $N = 16{,}384$. Per-level variance explained:
$\{0.9747, 0.9830, 0.9889\}$. The nested-subset invariant
(features at $k_i$ are a strict subset of features at $k_{i+1}$) holds
for $100\%$ of evaluated tokens by construction.

Compared to the flat TopK reference at $k = 64$: the Matryoshka at the
$k = 64$ level has slightly worse reconstruction ($0.983$ vs.\ $0.984$
variance explained) but dramatically fewer dead features ($1.64\%$ vs.\
$5.65\%$). The multi-level training appears to keep more capacity alive.

Decoder column cosine similarity across architectures:

\begin{center}
\begin{tabular}{lcc}
\toprule
Threshold & Count above & Fraction \\
\midrule
cos $> 0.50$ & $12{,}781$ & $0.780$ \\
cos $> 0.70$ & $9{,}517$ & $0.581$ \\
cos $> 0.90$ & $1{,}745$ & $\mathbf{0.106}$ \\
cos $> 0.95$ & $144$ & $0.009$ \\
cos $> 0.99$ & $4$ & $0.0002$ \\
\bottomrule
\end{tabular}
\end{center}

\textbf{Cross-architecture stability is even worse than cross-seed:}
$10.6\%$ of features at cos $> 0.9$, vs.\ $19.9\%$ cross-seed within the
same architecture. Matryoshka and flat TopK find substantively different
feature bases on the same model and the same data.

\subsection{Joint stability across all conditions}
\label{sec:navigator}

For each of the $16{,}384$ features in the reference SAE, we count the
number of cross-validation conditions in which it has a best-match cosine
$> 0.9$. The conditions are: $4\text{k}$-width SAE, $64\text{k}$-width
SAE, seed-$42$ SAE, Matryoshka SAE. Maximum possible score: $4$.

\begin{table}[h]
\centering
\caption{Joint stability across all four cross-validation conditions.
Only $2.14\%$ of features survive every comparison at cos $> 0.9$.}
\label{tab:joint}
\begin{tabular}{lrr}
\toprule
Stability score (cos $> 0.9$) & \# features & Fraction \\
\midrule
$0 / 4$ & $11{,}974$ & $73.08\%$ \\
$1 / 4$ & $1{,}799$ & $10.98\%$ \\
$2 / 4$ & $1{,}272$ & $7.76\%$ \\
$3 / 4$ & $988$ & $6.03\%$ \\
$\mathbf{4 / 4}$ & $\mathbf{351}$ & $\mathbf{2.14\%}$ \\
\bottomrule
\end{tabular}
\end{table}

The $351$ fully stable features are the strongest ``real feature''
candidates in this setup --- their decoder direction replicates regardless
of width, seed, or architectural objective.

\textbf{Of the $26$ auto-labeled monosemantic features, only $3$ are
fully stable:}
\begin{itemize}
\item f12117 --- BibTeX/LaTeX citation references (cos $\geq 0.997$
across every comparison)
\item f5747 --- File paths and directory structures (cos $\geq 0.997$
across every comparison)
\item f12697 --- Logical operators and negation symbols (cos $\geq 0.977$
across every comparison)
\end{itemize}
The remaining $23$ labeled features have stability score $\leq 1/4$:
they are clearly identified by auto-interp but do not survive any
non-trivial replication test. We take this as direct evidence that
\textbf{auto-interpretation overstates feature reliability by a factor of
approximately $8\times$} on this setup ($88.5\%$ monosemantic vs.\
$11.5\%$ fully stable).

\subsection{Multi-metric scorecard}
\label{sec:scorecard}

The six axes evaluated on the reference SAE:

\begin{center}
\begin{tabular}{lcl}
\toprule
Axis & Value & Source \\
\midrule
Reconstruction (var.\ explained) & $\mathbf{0.984}$ & basic eval \\
Predictive preservation (CE-delta loss recovered) & $\mathbf{0.954}$ & \S\ref{sec:ce-delta} \\
Capacity utilisation ($1 -$ dead fraction) & $\mathbf{0.948}$ & basic eval \\
Auto-interp monosemanticity rate & $0.885$ & \S\ref{sec:autointerp} \\
Cross-seed stability (cos $> 0.9$) & $\mathbf{0.199}$ & \S\ref{sec:seed} \\
Cross-architecture stability (cos $> 0.9$) & $\mathbf{0.106}$ & \S\ref{sec:matryoshka} \\
\bottomrule
\end{tabular}
\end{center}

The range is $0.106$ to $0.984$, a $\sim 88$ percentage-point gap. The
six axes do not agree. Reporting only the strongest axis (variance
explained) overstates the SAE's quality by nearly an order of magnitude
relative to the weakest axis (cross-architecture stability).

\section{Discussion}

\subsection{The standard suite is necessary but not sufficient}

Reconstruction, sparsity, and monosemanticity-rate are useful but
incomplete. They measure how well the SAE can encode and decode a given
model's activations and how interpretable the recovered features look.
They do not measure whether the features are a property of the model or
an artifact of the SAE training run. Our cross-validation results
indicate that, on this setup, the standard suite overstates feature
reliability by $\sim 8\times$.

The implication is not that SAEs are useless --- the $3$ fully stable
features we recover are clean, meaningful, and would survive any
reasonable interpretability test. The implication is that without
replication-style metrics, an SAE quality report cannot be cited as
evidence about the underlying model.

\subsection{What is the field's current practice?}

Most published SAE work reports reconstruction-style metrics on a single
trained SAE. A growing minority reports CE-delta substitution
\cite{lieberum2024gemma}. Cross-seed replication is occasionally
acknowledged as a limitation \cite{templeton2024scaling,
rajamanoharan2024gated} but rarely measured quantitatively at the
feature level. Cross-architecture replication (different SAE objectives
on the same model) is, to our knowledge, not standard practice.

We argue that cross-seed replication is the cheapest improvement
available --- one extra training run, identical hyperparameters --- and
should be reported alongside reconstruction in any SAE quality claim.

\subsection{The shortlist of stable features as a basis for downstream work}

The $351$ fully stable features ($2.14\%$ of the reference SAE) form a
small, well-validated candidate set for downstream interpretability
work --- circuit discovery, causal interventions, agent-behavior
auditing. Building circuits on the full SAE risks finding patterns that
are themselves training-noise artifacts.

\subsection{Tools we release}

\begin{enumerate}
\item \texttt{11\_sae\_scorecard.py} --- given a trained SAE checkpoint
and the standard set of evaluation JSON outputs, prints the six-axis
quality profile (\S\ref{sec:scorecard}).
\item \texttt{12\_feature\_navigator.py} --- given a reference SAE and
any set of comparison SAEs, computes per-feature stability scores and
emits both a JSON for programmatic use and a Markdown shortlist of
fully stable features. Supports \texttt{-{}-query <feature\_id>} for
single-feature lookup.
\end{enumerate}

\section{Limitations}

\begin{itemize}
\item \textbf{Single layer, single model.} We evaluate one residual stream
at one layer of one $160$M-parameter model. The numerical claims of
$2.14\%$ joint stability or $8\times$ overstatement should not be
read as universal. The methodological claim --- that
reconstruction-style and replication-style metrics can disagree by an
order of magnitude --- is more transportable.
\item \textbf{Single random seed pair.} Cross-seed stability is measured
across exactly two SAEs (reference and seed-$42$). Stability across
more seeds may yield different (likely lower) numbers.
\item \textbf{No causal validation.} We do not perform feature ablations
to test that fully stable features are also causally load-bearing for
downstream behavior. Cross-validation stability is a cheap proxy, not
a substitute for causal validation \cite{conmy2024automated,
marks2024feature}.
\item \textbf{Auto-interpretation is partial.} The Kimi-K2 labeling
collapses fine-grained splits we know exist (e.g., $11$ newline-variant
features all labeled ``newline tokens in various contexts''). A stronger
labeling model would likely raise the monosemanticity rate further
without changing the cross-validation stability, which would widen the
reconstruction-vs.-replication gap.
\item \textbf{Computational match.} A linear-encoder TopK SAE is
computationally matched to a single MLP layer; this constrains the
class of features it can find. Stronger SAE variants (multi-layer
encoders, attention-based encoders) could fit features the model itself
cannot use, exacerbating the cross-validation gap.
\item \textbf{$1$M tokens is small.} Modern frontier-scale SAE work uses
billions of activations. Some quantitative findings (e.g., the
$N \approx 4\text{k}$ effective capacity, the $63\%$ dead-feature rate
at $N = 65{,}536$) may shift at larger training scales.
\end{itemize}

\section{Conclusion}

We evaluated a TopK SAE on Pythia-160M layer $6$ along six independent
quality axes including three cross-validation tests. The standard
reconstruction-style metrics (variance explained, CE-delta loss
recovered) and the standard interpretability metric (auto-interp
monosemanticity rate) all score the SAE highly --- between $0.88$ and
$0.98$. The replication-style metrics (cross-seed, cross-architecture,
cross-width consistency) score the same SAE between $0.10$ and $0.20$.
The gap is large enough that the standard suite cannot be the only
quality claim attached to an SAE: it overstates feature reliability by
roughly $8\times$ on this setup.

We release two tools --- a multi-metric scorecard and a per-feature
stability navigator --- that surface this gap on any SAE. We expect
cross-seed replication to become standard in SAE reports; the marginal
cost is one extra training run, and the marginal information is
substantial.

\section*{Code and Artifacts}

\url{https://github.com/OE-GOD/sae-pythia-160m}

All checkpoints regenerable via the scripts in \texttt{code/}; raw
analysis outputs in \texttt{data/}; the trained $16$k reference SAE,
$4$k and $64$k width SAEs, seed-$42$ SAE, and Matryoshka SAE are
described and reproducible via the README.

\section*{Acknowledgments}

This work was conducted as a solo project as part of a series of
small-scale mechanistic interpretability replications. The author thanks
the Anthropic, OpenAI, and DeepMind interpretability teams for publishing
the methodology and results that made this small-scale replication
possible.

\begin{thebibliography}{99}

\bibitem{bricken2023monosemanticity}
Bricken, T., Templeton, A., et al. Towards Monosemanticity: Decomposing
Language Models With Dictionary Learning. \emph{Anthropic Transformer
Circuits Thread}, 2023.
\url{https://transformer-circuits.pub/2023/monosemantic-features}

\bibitem{templeton2024scaling}
Templeton, A. et al. Scaling Monosemanticity: Extracting Interpretable
Features from Claude 3 Sonnet. \emph{Anthropic Transformer Circuits
Thread}, 2024.
\url{https://transformer-circuits.pub/2024/scaling-monosemanticity}

\bibitem{gao2024topk}
Gao, L., Tour, T.D., Tillman, H., Goh, G., Troll, R., Radford, A.,
Sutskever, I., Leike, J., Wu, J. Scaling and Evaluating Sparse
Autoencoders. \emph{OpenAI}, 2024.
\url{https://arxiv.org/abs/2406.04093}

\bibitem{rajamanoharan2024gated}
Rajamanoharan, S., Conmy, A., Smith, L., Lieberum, T., Varma, V.,
Kram\'ar, J., Shah, R., Nanda, N. Improving Dictionary Learning with
Gated Sparse Autoencoders. \emph{DeepMind}, 2024.
\url{https://arxiv.org/abs/2404.16014}

\bibitem{rajamanoharan2024jumprelu}
Rajamanoharan, S., Lieberum, T., Sonnerat, N., Conmy, A., Varma, V.,
Kram\'ar, J., Nanda, N. Jumping Ahead: Improving Reconstruction Fidelity
with JumpReLU SAEs. \emph{DeepMind}, 2024.
\url{https://arxiv.org/abs/2407.14435}

\bibitem{lieberum2024gemma}
Lieberum, T., Rajamanoharan, S., Conmy, A., et al. Gemma Scope: Open
Sparse Autoencoders Everywhere All At Once on Gemma 2. \emph{DeepMind},
2024.
\url{https://arxiv.org/abs/2408.05147}

\bibitem{bussmann2025matryoshka}
Bussmann, B. et al. Learning Multi-Level Features with Matryoshka Sparse
Autoencoders. 2025.

\bibitem{biderman2023pythia}
Biderman, S., Schoelkopf, H., Anthony, Q.G., et al. Pythia: A Suite for
Analyzing Large Language Models Across Training and Scaling. \emph{ICML},
2023. \url{https://arxiv.org/abs/2304.01373}

\bibitem{conmy2024automated}
Conmy, A., Mavor-Parker, A., Lynch, A., Heimersheim, S., Garriga-Alonso,
A. Towards Automated Circuit Discovery for Mechanistic Interpretability.
\emph{NeurIPS}, 2023. \url{https://arxiv.org/abs/2304.14997}

\bibitem{marks2024feature}
Marks, S., Rager, C., Michaud, E.J., Belinkov, Y., Bau, D., Mueller, A.
Sparse Feature Circuits: Discovering and Editing Interpretable Causal
Graphs in Language Models. \emph{ICLR}, 2025.
\url{https://arxiv.org/abs/2403.19647}

\end{thebibliography}

\appendix

\section{The L1 Shrinkage Failure Mode We Started With}
\label{app:l1-shrinkage}

Before switching to TopK, we trained a vanilla $L_1$ SAE on the same
setup and swept $\lambda \in \{0.005, 0.02, 0.1\}$ at $N = 16{,}384$,
$1500$ training steps each. Despite a $20\times$ increase in
sparsity-penalty strength, the average $L_0$ remained between $1{,}048$
and $1{,}333$ active features per token --- well above the target band
of $30$-$100$.

\begin{table}[h]
\centering
\caption{$L_1$ SAE sparsity sweep.}
\label{tab:l1-sweep}
\begin{tabular}{lccc}
\toprule
$\lambda$ & Final MSE & Var.\ explained & $L_0$ \\
\midrule
$0.005$ & $2.80$ & $0.998$ & $1333$ \\
$0.02$  & $5.20$ & $0.997$ & $1302$ \\
$0.1$   & $14.72$ & $0.991$ & $1048$ \\
\bottomrule
\end{tabular}
\end{table}

This is the documented ``$L_1$ shrinkage'' failure mode
\cite{rajamanoharan2024gated}: minimising $\lambda \sum_i |f_i|$ is
indifferent between ``one feature firing at strength $S$'' and ``$S$
features firing at strength $1$''. Switching to TopK fixes this by hard
construction; $L_0 = k$ becomes exact.

\section{Selected Auto-Labeled Features}
\label{app:features}

The full $26$-feature catalog is in
\texttt{data/feature\_catalog.json}. Representative monosemantic
features:

\begin{tabular}{rl}
\toprule
ID & Auto-interpretation \\
\midrule
f1989 & Decimal points in numerical contexts \\
f5196 & Exponentiation notation in mathematical contexts \\
f5747 & Tokens related to file paths and directory structures \\
f7448 & Newline tokens inside code blocks \\
f10045 & Subword BPE continuations of multi-piece words \\
f12117 & BibTeX/LaTeX citation references \\
f12697 & Logical operators and negation symbols \\
\bottomrule
\end{tabular}

\textbf{Fully stable across cross-validation (cos $> 0.9$ in all 4
comparisons):} f12117, f5747, f12697.

\end{document}
